As observability becomes an increasingly indispensable property of microservice applications, and configuration becomes increasingly complex, there is also a growing concern regarding how to best configure and instrument an application. To enable developers and practitioners to assess the suitability of different designs, we set out to make observability a testable and quantifiable system property, similar to software quality measures like test coverage.  

In this paper, we presented an approach for assessing and improving the fault observability of cloud-based microservice applications. Our approach consists of a model for understanding and documenting the observability design space, a set of metrics to assess fault observability, and a tool that provides quantitative evidence for observability design decisions beyond gut-feeling or professional intuition.
We showed how this approach can be used in practice, with an application of the model to document the observability design space of a state-of-the-art cloud-based microservice application. Further, we evaluated multiple designs, revealing the until now unexplored observability-cost trade-offs.
We aim to improve \name{} by adding support for more platforms such as Kubernetes, by integrating the tool in CI pipelines, and by increasing the automation and intelligence of the assessment process.
With the help of other practitioners, the model and tooling can also be extended to cover more fault scenarios and more diverse microservice environments. With this work, we lay the foundation toward a systematic method for supporting observability design decisions in cloud-native microservice applications.

